% 自动生成文件：由 工程文件/tools/gen_appendix.py 生成，请勿手工编辑
\section{译注汇总}
\label{app:notes}
\noindent 本附录按节重列正文中的全部译者注，共 \textbf{57} 条（与正文脚注一一对应）. 编号「节号-序号」只用于本附录检索；来源等级与访问日期见各条末尾括注.
\begingroup\small
\medskip\noindent\textbf{第 2.1 节\quad 预备知识}\par\nopagebreak
\begin{description}[leftmargin=5.8em,labelwidth=5em,labelsep=0.6em,itemsep=0.15ex,topsep=0.2ex]
\item[2.1-1] 这里的 $2+2$ 指四维时空分解为紧致化所用的黎曼曲面 $C$（2 维）与剩余二维流形 $\Sigma$；$d$ 维理论对 $d-2$ 维流形赋予范畴，故要得到与 $C$ 有关的范畴必须取 $d=4$，同时剩余方向恰为二维，才能约化出 $\Sigma$ 上的二维 $\sigma$-模型. 原文 §2.1.（来源：A 级，arXiv:1506.04293 §2.1；A 级，arXiv:hep-th/0604151（Kapustin–Witten 2006）摘要，2026-11-19 核对.）
\item[2.1-2] 所谓“部分”拓扑场论，指配分函数一般并不良定义：四维流形 $M$ 上经典真空的模空间是模去共轭的 $\pi_1(M)\to G_{\mathbb C}$ 的同态空间，一般非紧，故路径积分与切割、粘合性质都难以证明；但膜范畴等不涉及该非紧性的可观测量仍良定义，这正是本篇只需膜范畴即可推进的原因. 原文 §2.1 作者原注.（来源：A 级，arXiv:1506.04293 §2.1 原注；A 级，arXiv:hep-th/0604151，2026-11-19 核对.）
\item[2.1-3] “希格斯丛”借自粒子物理而含义不同：这里的希格斯场 $\varphi$ 是 $C$ 上取值于 $\ade$ 的一次形式，并非标量场；它来自四维 $\mathcal N=4$ 理论经拓扑扭曲后沿 $C$ 约化所保留的数据. 该模空间带 $I,J,K$ 三个复结构，故称超凯勒流形. 原文 §2.1.（来源：B 级 nLab “Higgs bundle” \url{https://ncatlab.org/nlab/show/Higgs+bundle}，2026-11-19 访问；与 A 级 arXiv:1506.04293 §2.1–§2.2 交叉核对.）
\end{description}
\medskip\noindent\textbf{第 2.2 节\quad 希钦方程}\par\nopagebreak
\begin{description}[leftmargin=5.8em,labelwidth=5em,labelsep=0.6em,itemsep=0.15ex,topsep=0.2ex]
\item[2.2-1] 这组方程由 Hitchin 在 1987 年研究黎曼曲面上的自对偶方程时引入；把 $A$ 与 $\phi$ 合并为复值联络后，条件 (2.2) 就是四维自对偶方程沿 $C$ 约化的形式. 注意 $\phi$ 取值于 $\ade$ 且是一次形式，这是它与粒子物理中标量希格斯场的本质差别. 原文 §2.2 式 (2.1)–(2.2).（来源：A 级，Hitchin, Proc. London Math. Soc. (3) 55 (1987) 59–126（本文文献 [13]）；B 级 nLab “Higgs bundle”，2026-11-19 访问.）
\item[2.2-2] 复结构 $J$ 来自“半稳定同态 $\pi_1(C)\to G_{\mathbb C}$ 的模空间”这一描述，复结构 $I$ 则来自希格斯丛；两种描述之所以一致，依据的是非阿贝尔 Hodge 对应，而本文所用的半稳定性条件正是它的前提. 原文 §2.2.（来源：A 级，arXiv:1506.04293 §2.2；B 级 nLab “Higgs bundle” 中 nonabelian Hodge theory 一节，2026-11-19 访问.）
\item[2.2-3] $I,J,K$ 满足四元数关系 $IJ=K$、$I^2=J^2=K^2=-1$，故在每点切空间上给出四元数模结构；超凯勒流形上还有以单位球面参数化的复结构族，本篇只用到其中的 $I$ 与 $J$. 原文 §2.2.（来源：A 级，arXiv:1506.04293 §2.2；A 级，Hitchin, Proc. London Math. Soc. 55 (1987)（本文文献 [13]），2026-11-19 核对.）
\item[2.2-4] 超凯勒商要求两个矩映射同时为零：实矩映射对应 $D\star\phi=0$，复矩映射对应 $F-\phi\wedge\phi=0$，合起来正是式 (2.1) 的两组方程；文中“超凯勒矩映射为零”即指此. 原文 §2.2 式 (2.1)–(2.4).（来源：A 级，arXiv:1506.04293 §2.2；A 级，arXiv:hep-th/0604151，2026-11-19 核对.）
\item[2.2-5] 几何朗兰兹所用的 $A$-模型取辛形式 $\omega_K$ 而非 $\omega_J$，这是与电磁对偶相配的选择：在 $S$-对偶下 $A$-模型与 $B$-模型互换、$G$ 与 $^L\neg G$ 互换，对应关系即由此得到. 原文 §2.2.（来源：A 级，arXiv:1506.04293 §2.2；B 级 nLab “geometric Langlands correspondence” \url{https://ncatlab.org/nlab/show/geometric+Langlands+correspondence}，2026-11-19 访问；A 级，arXiv:hep-th/0604151.）
\end{description}
\medskip\noindent\textbf{第 2.3 节\quad $\MH$ 与余切丛}\par\nopagebreak
\begin{description}[leftmargin=5.8em,labelwidth=5em,labelsep=0.6em,itemsep=0.15ex,topsep=0.2ex]
\item[2.3-1] $T^*\mathcal M$ 的像在 $\MH$ 中是稠密开子集：$E$ 不稳定的稳定希钦对虽存在，但构成高余维集合，去掉后许多不变量不变，这正是可以用 $T^*\mathcal M$ 近似 $\MH$ 的原因；文献 [1] 由此把 $\MH$ 上的 $K$ 型 $A$-膜与 $\mathcal M$ 上的 $\mathcal D$-模联系起来. 原文 §2.3.（来源：A 级，arXiv:1506.04293 §2.3；A 级，arXiv:hep-th/0604151，2026-11-19 核对.）
\end{description}
\medskip\noindent\textbf{第 2.4 节\quad 希钦纤维化}\par\nopagebreak
\begin{description}[leftmargin=5.8em,labelwidth=5em,labelsep=0.6em,itemsep=0.15ex,topsep=0.2ex]
\item[2.4-1] 这里的 $2^{3g-3}$ 是 Bézout 式的启发计数：$3g-3$ 个二次方程一般有 $2^{3g-3}$ 个解，它只说明 $\ppi$ 在一般点处有限、从而满射，并非纤维的真实点数；由 §2.6 的谱曲线与阿贝尔化描述，纤维实为半维数的复环面. 原文 §2.4.（来源：A 级，arXiv:1506.04293 §2.4；B 级 nLab “Hitchin fibration” \url{https://ncatlab.org/nlab/show/Hitchin+fibration}，2026-11-19 访问.）
\end{description}
\medskip\noindent\textbf{第 2.5 节\quad 完全可积性}\par\nopagebreak
\begin{description}[leftmargin=5.8em,labelwidth=5em,labelsep=0.6em,itemsep=0.15ex,topsep=0.2ex]
\item[2.5-1] $\alpha_{a,i}$ 的取值可由塞尔对偶性算出：$\mathcal V$ 的切空间为 $\oplus_i H^0(C,K^{d_i})$，其对偶为 $\oplus_i H^1(C,K^{1-d_i})$，因为 $K\otimes K^{-d_i}=K^{1-d_i}$；故 $\alpha_{a,i}\in H^1(C,K^{1-d_i})$ 正是与 $H^0(C,K^{d_i})$ 配对的上同调类，$H_{a,i}$ 是底空间上的线性函数. 原文 §2.5 式 (2.7).（来源：A 级，arXiv:1506.04293 §2.5；塞尔对偶性的直接计算，2026-11-19 核对.）
\end{description}
\medskip\noindent\textbf{第 2.6.1 节\quad 谱曲线}\par\nopagebreak
\begin{description}[leftmargin=5.8em,labelwidth=5em,labelsep=0.6em,itemsep=0.15ex,topsep=0.2ex]
\item[2.6.1-1] 可用黎曼–胡尔维茨公式复核：$2g_D-2=N(2g_C-2)+\deg\mathcal R$，而 $\deg\mathcal R=2N(N-1)(g_C-1)$，故 $g_D=N^2(g_C-1)+1=g_C+(N^2-1)(g_C-1)$；判别式的次数 $N(N-1)$ 正是 $\mathcal R$ 的像在 $C$ 上的次数. 原文 §2.6.1.（来源：A 级，arXiv:1506.04293 §2.6.1；A 级，Hitchin, Duke Math. J. 54 (1987) 91–114（本文文献 [21]），2026-11-19 核对.）
\item[2.6.1-2] ${\cal N}(\mathcal R)$ 是允许沿分歧除子 $\mathcal R$ 取极点的扭化丛，即 ${\cal N}\otimes\mathcal O_D(\mathcal R)$；它与式 (2.22) 的写法一致：$y$ 是 $\psi^*(K)$ 的截面、$\PP'(y)$ 关于 $y$ 为 $N-1$ 次，零点除子恰为 $\mathcal R$，于是 $\mathcal R$ 对应 $\psi^*(K^{N-1})$. 原文 §2.6.1 式 (2.21)–(2.22).（来源：A 级，arXiv:1506.04293 §2.6.1；A 级，Hitchin, Duke Math. J. 54 (1987)（本文文献 [21]），2026-11-19 核对.）
\end{description}
\medskip\noindent\textbf{第 2.6.2 节\quad 谱曲线}\par\nopagebreak
\begin{description}[leftmargin=5.8em,labelwidth=5em,labelsep=0.6em,itemsep=0.15ex,topsep=0.2ex]
\item[2.6.2-1] 这正是本节标题所说的“阿贝尔化”：把 $C$ 上非阿贝尔的 $G$-丛 $E$ 换成谱覆盖 $D$ 上的线丛 ${\cal L}$，代价是工作曲线从 $C$ 换成 $D$——$D\to C$ 是 $N$ 重覆盖，亏格为 $N^2(g_C-1)+1$，远大于 $g_C$；希格斯场则由 $\varphi=\psi_*(y)$ 一并恢复. 原文 §2.6.2 式 (2.25)、(2.29).（来源：A 级，arXiv:1506.04293 §2.6.2；A 级，arXiv:hep-th/0604151，2026-11-19 核对.）
\end{description}
\medskip\noindent\textbf{第 2.6.3 节\quad 谱曲线}\par\nopagebreak
\begin{description}[leftmargin=5.8em,labelwidth=5em,labelsep=0.6em,itemsep=0.15ex,topsep=0.2ex]
\item[2.6.3-1] 取基点后复环面才成为阿贝尔簇：$\Lambda({\cal L})=\det\psi_*({\cal L})$ 把 $\mathrm{Pic}_{c_0}(D)$ 映到 $\mathrm{Jac}(C)$，其核 $\mathcal F_{SU(N)}$ 的维数为 $g_D-g_C=\dim\mathcal F$，正是一个“Prym 型”对象（推前为平凡线丛的线丛构成的空间）. ${\cal L}_0=\psi^*(K^{(N-1)/2})$ 给出核中一个自然点：奇 $N$ 时它与 $K^{1/2}$ 的选取无关，偶 $N$ 时只差一个二阶平坦线丛的自由度，这一差别由文献 [23] 指出并推广到其他 $G$. 原文 §2.6.3.（来源：A 级，arXiv:1506.04293 §2.6.3；B 级 nLab “Prym variety” \url{https://ncatlab.org/nlab/show/Prym+variety}，2026-11-19 访问；A 级，arXiv:math/0005132（本文文献 [23]）.）
\end{description}
\medskip\noindent\textbf{第 2.6.4 节\quad }\par\nopagebreak
\begin{description}[leftmargin=5.8em,labelwidth=5em,labelsep=0.6em,itemsep=0.15ex,topsep=0.2ex]
\item[2.6.4-1] 这里“平坦 ${\mathrm{Spin}_c}$ 结构”的来历是 D-膜的 $K$ 理论分类：包裹子流形 $C$ 的 D-膜，其陈-帕顿丛必须与 $C$ 在 $W$ 中法丛的旋量丛张量在一起才有意义，而 $W=T^*C$ 的法丛同构于余切丛 $K$，故要求 $\U^2\cong K$。当 $w_2(C)\neq0$ 时 $K$ 没有整体平方根，只能存在平坦 ${\mathrm{Spin}_c}$ 结构，其次数为 $g_C-1$。（来源：A 级，Witten，\emph{D-Branes And K-Theory}，arXiv:hep-th/9810188 §4.3“The Spin\_c Case”，判据 $c_1(L)\equiv w_2(N)\bmod 2$；原文 §2.6.4 与式 (\ref{tywo})，2026-11-19 核对.）
\end{description}
\medskip\noindent\textbf{第 2.6.5 节\quad }\par\nopagebreak
\begin{description}[leftmargin=5.8em,labelwidth=5em,labelsep=0.6em,itemsep=0.15ex,topsep=0.2ex]
\item[2.6.5-1] “典范地成为阿贝尔簇”指这个复环面有一个不依赖额外选择的零点，从而可等同于 $\mathrm{Pic}^0$ 型的对象。$N$ 为奇数时，典范点 $\psi^*(K^{(N-1)/2})$ 由 $K$ 的整数次幂直接给出；$N$ 为偶数时 $(N-1)/2$ 不是整数，必须先取定 $K^{1/2}$（自旋结构），典范点才存在，这正是“还不能完全算作阿贝尔簇”的意思。后文 (2.38) 中除以 $\Gamma_N$ 恰好消去对自旋结构的依赖。（来源：A 级，原文 §2.6.3 与 §2.6.5 式 (2.36)–(2.38)；A 级，Frenkel–Witten，arXiv:0710.5939 §10 对相关修正的讨论；2026-11-19 核对.）
\end{description}
\medskip\noindent\textbf{第 2.6.6 节\quad }\par\nopagebreak
\begin{description}[leftmargin=5.8em,labelwidth=5em,labelsep=0.6em,itemsep=0.15ex,topsep=0.2ex]
\item[2.6.6-1] 这类特征类就是所谓 \emph{'t Hooft 磁通}：$C$ 是二维的，$PSU(N)=SU(N)/\Bbb Z_N$ 丛的分类只依赖 $\pi_1(PSU(N))=\Bbb Z_N$，故 $H^2(C,\Bbb Z_N)\cong\Bbb Z_N$；$N=2$ 时 $PSU(2)=SO(3)$，$\xi$ 就是 $w_2(E)$。若 $PSU(N)$ 丛来自 $U(N)$ 丛 $E$，则 $\xi\equiv c_1(E)\bmod N$，这解释了下一段的提升判据。（来源：A 级，[KW]（arXiv:hep-th/0604151）§7.1 关于 $G_{ad}$ 丛特征类 $\xi(E)\in H^2(M,\Bbb Z)$ 与 $SO(3)$ 情形即 $w_2$ 的讨论；原文 §2.6.6；2026-11-19 核对.）
\item[2.6.6-2] 这里“非平凡的 gerbe”可理解为 $\sigma$-模型的离散 $B$-场：$U(1)$-gerbe 的同构类由 $H^2(X,U(1))$ 分类，其上的向量丛只能“扭曲地”存在，即转移函数只须满足 $v_{ki}v_{jk}v_{ij}=b_{ijk}$。$PSU(N)$ 分支与 $SU(N)$ 分支对偶时，两边相差的正是这份离散数据。（来源：A 级，[KW]（arXiv:hep-th/0604151）§7.2 关于平坦 $B$-场 $b\in H^2(X,U(1))$ 与 twisted vector bundle 的讨论；B 级，nLab“gerbe”\url{https://ncatlab.org/nlab/show/gerbe}（Last revised 2026-02-14，2026-11-19 访问）给出 gerbe 的叠定义，待核；原文 §2.6.6 与 §3.2.）
\end{description}
\medskip\noindent\textbf{第 2.6.7 节\quad }\par\nopagebreak
\begin{description}[leftmargin=5.8em,labelwidth=5em,labelsep=0.6em,itemsep=0.15ex,topsep=0.2ex]
\item[2.6.7-1] 对一般 $G$，“谱覆盖”要换成 \emph{cameral 覆盖}：它是 $C$ 的一个 $|W|$ 重覆盖（$W$ 为 Weyl 群），在正则点处的纤维是 $G$ 的 Cartan 子代数。谱覆盖只对 $GL(N)$ 或具有“小”表示的群才有简单意义，cameral 覆盖则对所有 $G$ 都有定义；两者的关系是：谱覆盖为 cameral 覆盖与小表示所做的纤维积。（来源：A 级，Donagi–Gaitsgory，\emph{The Gerbe Of Higgs Bundles}，arXiv:math/0005132 §2“Higgs bundles and cameral covers”与 §9“Spectral covers versus cameral covers for $G=GL(n)$”；原文 §2.6.7；2026-11-19 核对.）
\end{description}
\medskip\noindent\textbf{第 2.7.1 节\quad 特选截面}\par\nopagebreak
\begin{description}[leftmargin=5.8em,labelwidth=5em,labelsep=0.6em,itemsep=0.15ex,topsep=0.2ex]
\item[2.7.1-1] “特选截面”指每个希钦纤维里都有一个由不稳定丛给出的标准点：$SL(2,\Bbb C)$ 情形取 $E=K^{1/2}\oplus K^{-1/2}$，其上的希格斯场在 $w\neq0$ 时稳定，于是纤维化得到一个不依赖参数的标准截面；而对一般稳定 $E$，截面与纤维的交点数为 $2^{3g-3}$，说明这种“唯一标准点”只在退化（不稳定）丛处成立。（来源：A 级，Hitchin，\emph{The self-duality equations on a Riemann surface}，Proc. London Math. Soc. (3) 55 (1987) 59–126；原文 §2.7.1 与 §2.6.3；2026-11-19 核对.）
\end{description}
\medskip\noindent\textbf{第 2.7.2 节\quad 特选截面}\par\nopagebreak
\begin{description}[leftmargin=5.8em,labelwidth=5em,labelsep=0.6em,itemsep=0.15ex,topsep=0.2ex]
\item[2.7.2-1] “oper”是 Beilinson–Drinfeld 引入的概念：粗略地说，它是带有一个满足额外横截与正则性条件的联络的 $G$-丛（$\mathrm{PSL}(2)$ 情形下，oper 等同于 $C$ 上的射影结构）。oper 的模空间给出希钦纤维化的一个截面，其取值正是 Kostant 主 $SU(2)$ 嵌入所确定的不稳定丛，故截面不依赖纤维化参数，只因 $K^{1/2}$ 的选择而“几乎”典范。（来源：A 级，Beilinson–Drinfeld，\emph{Quantization Of Hitchin's Integrable System And Hecke Eigensheaves}，preprint，即原文 §2.7.2 所引 [12]，opers 概念与 Hitchin 系统的关系见该书导言；原文 §2.7.2；2026-11-19 核对.）
\end{description}
\medskip\noindent\textbf{第 3 节\quad 对偶环面与希钦纤维化}\par\nopagebreak
\begin{description}[leftmargin=5.8em,labelwidth=5em,labelsep=0.6em,itemsep=0.15ex,topsep=0.2ex]
\item[3-1] “唯一特选点”给出一种标准的主极化：把 $\cF$ 与 $\cF'$ 的格分别看作电荷格与磁荷格，固定 ${\cal T}$ 在唯一特选点处的平凡化，便可把 $\cF'$ 等同于 $\mathrm{Pic}^0(\cF)$。这正是 Goddard–Nuyts–Olive 对偶（GNO 对偶）的几何形式：一个群的电荷格与对偶群的磁荷格互为对偶。（来源：A 级，Goddard–Nuyts–Olive，\emph{Gauge Theories and Magnetic Charge}，Nucl. Phys. B125 (1977) 1–28；原文 §3 对偶环面定义；2026-11-19 核对.）
\end{description}
\medskip\noindent\textbf{第 3.1 节\quad 例子}\par\nopagebreak
\begin{description}[leftmargin=5.8em,labelwidth=5em,labelsep=0.6em,itemsep=0.15ex,topsep=0.2ex]
\item[3.1-1] 其机制是：陈-西蒙斯作用量 $I=(1/2\pi)\int A\wedge dB$ 在带边三维流形上的变分给出边界项 $\frac{1}{2\pi}\int_{\partial X}\delta A\wedge\delta B$，故 $\omega$ 正是 $I$ 的“边界辛形式”；$1/2\pi$ 的归一化使 $I$ 的周期属于 $2\pi\Bbb Z$，从而预量子化线丛 ${\cal T}$ 存在且在同构意义下唯一。（来源：A 级，Ramadas–Singer–Weitsman，\emph{Some Comments On Chern-Simons Gauge Theory}，Commun. Math. Phys. 126 (1989) 409；原文 §3.1 式 (3.1)–(3.3)；2026-11-19 核对.）
\end{description}
\medskip\noindent\textbf{第 3.3 节\quad 拓扑观点}\par\nopagebreak
\begin{description}[leftmargin=5.8em,labelwidth=5em,labelsep=0.6em,itemsep=0.15ex,topsep=0.2ex]
\item[3.3-1] 此处的 $\psi^*$ 与 §2.6.6、§3.2 中的 $\psi^*$ 同名而含义不同：(3.15) 的 $\psi^*:H_1(C)\to H_1(D)$ 把 $C$ 中的闭链提升为 $D$ 中的闭链，满足 $\psi_*\psi^*=N$（$N$ 为覆盖次数）；而线丛情形下的 $\psi^*:\mathrm{Pic}(C)\to\mathrm{Pic}(D)$ 是拉回。读数时须按上下文分辨。（来源：A 级，原文 §3.3 式 (3.15)(3.16) 与 §2.6.6 式 (2.36)–(2.39)；2026-11-19 核对.）
\item[3.3-2] 拓扑观点的要点是：$\FF_{SU(N)}=\mathrm{Hom}(\Gamma^*,U(1))$，只要 $\Gamma^*=H_1(D)/\psi^*H_1(C)$ 无挠，$\mathrm{Hom}(\Gamma^*,U(1))$ 就是紧环面，特别地连通；“无挠”来自 $\psi_*$ 满射，它保证 $\psi^*H_1(C)$ 是 $H_1(D)$ 的饱和子格。这与 §2.6.3 由复几何得到的预言互为独立证明。（来源：A 级，原文 §3.3 式 (3.19)–(3.23)；A 级，Hausel–Thaddeus，arXiv:math/0205236 对阿贝尔簇纤维的拓扑处理；2026-11-19 核对.）
\end{description}
\medskip\noindent\textbf{第 3.4 节\quad 辛形式}\par\nopagebreak
\begin{description}[leftmargin=5.8em,labelwidth=5em,labelsep=0.6em,itemsep=0.15ex,topsep=0.2ex]
\item[3.4-1] $\psi_*$ 是有限态射 $\psi:D\to C$ 的迹映射（在每根纤维上取迹），它满射，核的维数恰为 $g_D-g_C$。对 $SU(N)$ 有 $\Tr\varphi=0$，故 $\lambda=y\,dz$ 满足 $\psi_*(\lambda)=\sum_i y_i\,dz=0$，即 $\lambda$ 落在核中；于是 $\partial\lambda/\partial H_a$（$a=1,\dots,g_D-g_C$）给出核的一组基。（来源：A 级，原文 §3.4 式 (3.32)(3.33)；A 级，Hitchin，\emph{Stable bundles and integrable systems}，Duke Math. J. 54 (1987) 91–114 关于谱曲线与迹映射；2026-11-19 核对.）
\end{description}
\medskip\noindent\textbf{第 3.4.1 节\quad 辛形式}\par\nopagebreak
\begin{description}[leftmargin=5.8em,labelwidth=5em,labelsep=0.6em,itemsep=0.15ex,topsep=0.2ex]
\item[3.4.1-1] 这里的归一化是：约定 $|d^2z|=i\,dz\wedge d\bar z$，$\delta_{\rr_i}$ 满足 $\int_D\delta_{\rr_i}=1$；由 $\bar\partial\partial f=i\pi\sum_i n_i\delta_{\rr_i}$ 与 $f\sim\frac12\ln|z-\rr_i|^2$ 可核出 (3.38) 的系数 $-i\pi$（关键是 $\partial_{\rr_i}$ 作用在奇性上等于 $-\partial_z$），再乘 $|d^2z|$ 即得 (3.39) 与 (\ref{wellwrite}) 一致。原文该处脚注说“具有曲率 $f=da=2\pi i\sum_i n_i\delta_{\rr_i}$”，其中 $f$ 其实是标量函数而非曲率：由 $a_z=-\partial_zf$、$a_{\bar z}=\partial_{\bar z}f$ 得 $da=2\bar\partial\partial f=2\pi i\sum_i n_i\delta_{\rr_i}$，故曲率是 $2\pi i\sum_i n_i\delta_{\rr_i}$，与 (3.38) 的系数相差因子 2，两者并不矛盾。（来源：A 级，原文 §3.4.1 式 (3.37)–(3.39) 及其脚注；2026-11-19 核对.）
\end{description}
\medskip\noindent\textbf{第 4.1 节\quad 本征膜}\par\nopagebreak
\begin{description}[leftmargin=5.8em,labelwidth=5em,labelsep=0.6em,itemsep=0.15ex,topsep=0.2ex]
\item[4.1-1] 这里的 $V$ 只在同构意义下确定：若 $L\B\cong \B\otimes V$，则任何与 $V$ 同构的向量空间同样满足定义，且 $V$ 的自同构群 $GL(n,\C)$ 可用来构造 $V$ 在其中非平凡变化的膜族（上一句已提及）. 在二维拓扑场论中，膜构成一个 $\C$-线性范畴：对象是边界条件，态射是边界上的局域算符；线算符作用于膜给出该范畴的一个张量函子，因此“对角化”只能在范畴层面表述为同构 $L\B\cong \B\otimes V$，而不能写成数值等式，$V$ 的维数相当于 $L$ 在该膜上的“多重数”. 当 $\B$ 是支撑在点 $m$ 处的零维膜时，$V$ 由 $L$ 在 $m$ 处的纤维唯一决定，这正是零维膜作为本征膜时计算最简单的根本原因. [来源：nLab “electric eigenbrane”，\url{https://ncatlab.org/nlab/show/electric+eigenbrane}（Created 2014-07-02；2026-11-19 访问）；[KW]（arXiv:hep-th/0604151）§8（原文 §4.2 所引）；原文 §4.1 式 (4.3) 及上下文.]
\end{description}
\medskip\noindent\textbf{第 4.2 节\quad 电本征膜}\par\nopagebreak
\begin{description}[leftmargin=5.8em,labelwidth=5em,labelsep=0.6em,itemsep=0.15ex,topsep=0.2ex]
\item[4.2-1] 复结构 $J$ 的地位可以用 $\MH$ 的超凯勒三联 $(I,J,K)$ 说明：在 $I$ 中 $\MH$ 是希格斯丛的模空间，而在 $J$、$K$ 中它是 $C$ 上平坦 $^L\neg G$-丛（即局部系统）的模空间. 因此“$J$ 中的电本征膜”就是支撑在某个平坦丛处、其陈-帕顿丛带有平坦联络的膜；经超凯勒旋转，这类膜给出 $D$-模，即几何朗兰兹中的赫克本征层. 这正是原文“几何朗兰兹是关于复结构 $J$ 中的电本征膜的陈述”一语的技术含义，也是第 \ref{magbranes} 节要把电本征膜与磁本征膜纳入同一超凯勒框架的原因. [来源：nLab “Kapustin-Witten TQFT”，\url{https://ncatlab.org/nlab/show/Kapustin-Witten+TQFT}（Last revised 2026-09-28；2026-11-19 访问）；原文 §4.2、§5.1.]
\end{description}
\medskip\noindent\textbf{第 4.3 节\quad 't Hooft 算符与赫克变换}\par\nopagebreak
\begin{description}[leftmargin=5.8em,labelwidth=5em,labelsep=0.6em,itemsep=0.15ex,topsep=0.2ex]
\item[4.3-1] 式 (4.9) 的归一化可以显式验证：取 $S^2$ 为绕原点的球面，由 $\Delta(1/r)=-4\pi\delta^{(3)}$ 得 $\int_{S^2}F=\pm 2\pi i$（符号取决于定向），即磁通绝对值为 $2\pi$. 在复化约定下这等价于 $e^{\int_{S^2}F}=1$（若换回实场强，则写作 $e^{i\int_{S^2}F}=1$），即狄拉克量子化条件恰好被满足；常数 $i/2$ 就是为使荷为 1 的磁通恰好等于一个量子而选取的. 这也是 §4.4.1 中 $\O\to\O(m\pp)$ 的整数 $m$ 可以与磁荷等同的原因. [来源：nLab “monopole”，\url{https://ncatlab.org/nlab/show/monopole}（Last revised 2026-05-27；2026-11-19 访问；磁荷＝U(1) 主丛的第一陈类）；原文 §4.3 式 (4.9)、§4.4.1.]
\item[4.3-2] 这一分类在数学上称为 GNO（Goddard–Nuyts–Olive）对偶：$G$ 的磁荷（单极子荷）取值于 $\pi_1(G)$ 的对偶，而在共轭意义下，从 $U(1)$ 到 $G$ 的同态 $\rho$ 恰由一个优势余权刻画，对偶地由 $^L\neg G$ 的一个优势最高权 $\lambda$ 刻画；二者以 $\langle\lambda,\rho\rangle\in\Z$ 配对，这个整数就是 't Hooft 算符的荷. 因此 $G$ 的 't Hooft 算符的荷可以用 $^L\neg G$ 的最高权来标记，并与 $^L\neg G$ 的威尔逊算符一一对应. 原文只陈述这一对应，未展开配对的具体构造. [来源：nLab “’t Hooft operator”，\url{https://ncatlab.org/nlab/show/}\%27t+Hooft+operator（Last revised 2025-07-04；2026-11-19 访问）；[KW]（arXiv:hep-th/0604151）§6.2（该页所引）；原文 §4.3 式 (4.10).]
\item[4.3-3] 这里的两组写法是同一物理内容的两种形式：式 (4.11) 是 3 维欧氏空间上的博戈莫尔尼方程（即 4 维自对偶杨--米尔斯方程经平移不变约化到 3 维所得），而 $D_1,D_2,D_3$ 是超凯勒商构造中的标准记号——$D_1$ 是 $\bar\partial$ 型算符，$D_2,D_3$ 分别是另外两个复结构中的 $\bar\partial$ 型算符，$\mu=\sum_i[D_i,D_i^\dagger]$ 则是相应的超凯勒矩映射. 于是“复方程 $[D_i,D_j]=0$ 加 $\mu=0$”就是把希钦方程（从而把 $\MH$ 的复结构）写成三个复结构的统一形式，这也是原文说这些条件“在相当简单的应用中都满足”的原因. 另外，原文按 [1] 式 (3.29) 取 $t=1$，$t$ 是 [1] 中那一族拓扑扭曲理论的参数，取 $1$ 是 [1] 中为研究 A-模型而选取的方便值. [来源：原文 §4.3 式 (4.11)–(4.16) 及上下文；nLab “Bogomolny equation”，\url{https://ncatlab.org/nlab/show/Bogomolny+equation}（Created 2019-10-14；2026-11-19 访问）；[KW]（arXiv:hep-th/0604151）§3.6（原文所引）.]
\item[4.3-4] 在数学侧，$\pp$ 处的赫克修正的参数空间可以等同于仿射 Grassmannian $\mathrm{Gr}_G$ 的一个 Schubert 胞腔：$\mathrm{Gr}_G$ 的对象是形式圆盘 $D=\mathrm{Spec}\,\C[[z]]$ 上的 $G$-丛连同在穿孔圆盘 $D^*$ 上的平凡化，而平凡化在 $\pp$ 处允许的极点阶数由一个优势余权（对偶地，$^L\neg G$ 的一个优势最高权）标记，对应的 Schubert 胞腔就是这一类型的全部赫克修正. 因此“赫克修正的类型”天然由 $^L\neg G$ 的最高权分类；当该权微小、相应胞腔紧致时不需要单极子冒泡，这正是下一小节判据的几何来源. 原文只引 [1] 的 §9、§10，未给出仿射 Grassmannian 的构造. [来源：nLab “affine Grassmannian”，\url{https://ncatlab.org/nlab/show/affine+Grassmannian}（Last revised 2023-08-15；2026-11-19 访问）；[KW]（arXiv:hep-th/0604151）§9–§10（原文 §4.3 所引）；原文 §4.3.]
\end{description}
\medskip\noindent\textbf{第 4.4.1 节\quad 基本例子}\par\nopagebreak
\begin{description}[leftmargin=5.8em,labelwidth=5em,labelsep=0.6em,itemsep=0.15ex,topsep=0.2ex]
\item[4.4.1-1] 更定量地说，$\O(m\pp)$ 与 $\O$ 的差别由第一陈数给出：$c_1(\O(m\pp))=m$，而 $\int_C c_1=m$ 正是穿过 $C$ 的磁通数（以式 (4.9) 的 $2\pi i$ 为单位）. 线算符可以融合，$T(m_1)\cdot T(m_2)=T(m_1+m_2)$，对应地 $\O\to\O(m_1\pp)\to\O((m_1+m_2)\pp)$，磁荷相加；$m<0$ 时修正方向相反. 这给出了 $U(1)$ 情形“赫克修正的类型就是磁荷”的完整字典. [来源：nLab “monopole”，\url{https://ncatlab.org/nlab/show/monopole}（Last revised 2026-05-27；2026-11-19 访问；磁荷＝第一陈类）；[KW]（arXiv:hep-th/0604151）§4.5（原文 §4.4.1 所引）；原文 §4.4.1.]
\end{description}
\medskip\noindent\textbf{第 4.4.2 节\quad 基本例子}\par\nopagebreak
\begin{description}[leftmargin=5.8em,labelwidth=5em,labelsep=0.6em,itemsep=0.15ex,topsep=0.2ex]
\item[4.4.2-1] 这个 $\CP^1=\Bbb P(E_{-,\pp})$ 是仿射 Grassmannian $\mathrm{Gr}_{U(2)}$ 中对应于微小权 $(1,0)$ 的 Schubert 胞腔：$\pp$ 处的一阶极点配上留数方向 $(u,v)$，正好给出该胞腔的全部（有限维的）参数. 对 $U(N)$ 的微小权 $(1,\dots,1,0,\dots,0)$，同样的构造给出有限维的 $\mathbf{Gr}(k,N)$（见式 (4.30) 之后的讨论），因为此时 Schubert 胞腔紧致；一般权的 Schubert 胞腔不紧致，其紧化需要粘上较低权的胞腔，物理上就是下一段所说的单极子冒泡. [来源：nLab “affine Grassmannian”，\url{https://ncatlab.org/nlab/show/affine+Grassmannian}（Last revised 2023-08-15；2026-11-19 访问）；[KW]（arXiv:hep-th/0604151）§10（原文 §4.4.2 所引）；原文 §4.4.2.]
\item[4.4.2-2] 从物理上看，紧化的粘合点对应“在一个 't Hooft 算符的背景中，一个光滑 BPS 单极子收缩到零尺度”：单极子带磁荷 $k$ 时会把部分磁荷交给 't Hooft 算符，使原来的权 $(n,m)$ 变为较低的权（例如 $(n-k,m+k)$），这正是上一段所说的“较低权”修正；单极子半径趋于零时总磁荷守恒，所以紧化中只会出现同一总荷的较低权修正. 由于需要 $1\le k\le (n-m)/2$，只有 $n-m\ge 2$ 时才会发生冒泡，这解释了为何 $n-m\le 1$（微小情形）的修正空间本身就是紧的、一个 Schubert 胞腔就够了. [来源：[KW]（arXiv:hep-th/0604151）§10（原文 §4.4.2 所引）；nLab “monopole”，\url{https://ncatlab.org/nlab/show/monopole}（Last revised 2026-05-27；2026-11-19 访问）；原文 §4.4.2 式 (4.23) 及脚注.]
\item[4.4.2-3] 判断微小有一个实用判据：不可约表示 $^L\neg R$ 微小，当且仅当其维数等于最高权的 Weyl 群轨道大小，等价地，当且仅当其所有权同属一个 Weyl 轨道. 对 $U(2)$，权为 $(n-j,m+j)$（$j=0,1,\dots,n-m$），Weyl 群 $S_2$ 交换两个分量，所以恰好当 $n-m\le 1$ 时所有权在同一轨道上——这与原文的条件 $n-m\le1$ 一致. 特别地，$SU(2)$ 只有二维表示微小；$SO(3)$ 的中心平凡，故没有非平凡的微小表示（与“微小表示对应中心特征标”的一般事实相符）. [来源：原文 §4.4.2（微小表示的定义及式 (4.23) 之后的讨论）；[KW]（arXiv:hep-th/0604151）§10；2026-11-19 核对.]
\end{description}
\medskip\noindent\textbf{第 4.4.3 节\quad 基本例子}\par\nopagebreak
\begin{description}[leftmargin=5.8em,labelwidth=5em,labelsep=0.6em,itemsep=0.15ex,topsep=0.2ex]
\item[4.4.3-1] 用 $V=\ad_0(E)$ 而不直接用 $E$ 的必要性在于：把 $E$ 换成 $E\otimes\L$（$\L$ 为任意线丛）不改变 $\ad_0(E)$，因此 $E_\pm$ 只是 $SO(3)$ 丛 $V_\pm$ 的一种“提升”，只有 $V_\pm$ 才有典范意义；式 (4.26)、(4.27) 都应理解为式 (4.25) 的简写. 顺带指出：§4.4.3 开头写 $V_-=\ad(E_-)=\O\oplus\O\oplus\O$ 时，$\ad(E_-)$ 应是 $\ad_0(E_-)$——$\ad(E_-)=E_-\otimes E_-^*$ 是 4 维的，其无迹部分才是 3 维的 $\O\oplus\O\oplus\O$（原文 p.46 同此，属原文笔误）. [来源：原文 §4.4.3 式 (4.25)–(4.27) 及上下文（原文 p.46）；2026-11-19 核对.]
\end{description}
\medskip\noindent\textbf{第 4.5.2 节\quad 纳入希格斯场}\par\nopagebreak
\begin{description}[leftmargin=5.8em,labelwidth=5em,labelsep=0.6em,itemsep=0.15ex,topsep=0.2ex]
\item[4.5.2-1] $\varphi_-(\pp)$ 在 $\CP^1$ 上生成一个全纯向量场（在本征基下形如 $u\partial_u-v\partial_v$），其零点恰为两个不动点 $[1:0]$ 与 $[0:1]$，分别对应 $\varphi_-(\pp)$ 的两个本征方向. 因此“$\CP^1$ 中的点 $\bb$ 在 $\varphi_-(\pp)$ 生成的对称性下不变”就是要求 $\bb$ 是本征向量方向，而 $\CP^1$ 上恰有两个这样的点——这就是下一段“恰有两个 $\varphi_-(\pp)$-不变赫克修正”的缘由，也是它与谱覆盖的两个点 $\qq',\qq''$ 一一对应的几何解释. 对一般群，$\varphi_-(\pp)$ 的本征分解同样给出有限个不变点，其个数由 Weyl 群控制，§4.5.5 中的 $N!/\#\Gamma$ 即其显式形式. [来源：原文 §4.5.2 式 (4.41)–(4.43) 及上下文；原文 §4.5.5；2026-11-19 核对.]
\end{description}
\medskip\noindent\textbf{第 4.5.5 节\quad 纳入希格斯场}\par\nopagebreak
\begin{description}[leftmargin=5.8em,labelwidth=5em,labelsep=0.6em,itemsep=0.15ex,topsep=0.2ex]
\item[4.5.5-1] 这个数可以这样理解：$\varphi_-(\pp)$ 有 $N$ 个互异本征值，在 $\pp$ 处的赫克修正空间中，$\varphi_-(\pp)$-不变的点的个数等于“把 $N$ 个本征方向分配给 $N$ 个权 $m_1,\dots,m_N$”的方式数；当 $m_i$ 两两互异时有 $N!$ 种，而当 $m_i$ 出现重数时，重数相同的指标可以任意置换，故要除以稳定子群 $\Gamma$ 的阶. 在谱覆盖上，这就是把 $D$ 上位于 $\pp$ 上方的 $N$ 个点 $\qq_1,\dots,\qq_N$ 作各种排序；式 (4.49) 的 $\L\to \L\otimes(\otimes_i\O(\qq_i)^{m_i})$ 正是在说赫克修正只改变线丛的因子、不改变谱覆盖. 这也与 $U(2)$ 情形两个 $m$ 互异时恰好两个不变修正相一致. [来源：原文 §4.5.5 式 (4.47)–(4.49) 及上下文；nLab “Hitchin fibration”，\url{https://ncatlab.org/nlab/show/Hitchin+fibration}（Last revised 2023-11-01；2026-11-19 访问）；2026-11-19 核对.]
\end{description}
\medskip\noindent\textbf{第 5.2 节\quad 在希钦纤维化的纤维上的作用}\par\nopagebreak
\begin{description}[leftmargin=5.8em,labelwidth=5em,labelsep=0.6em,itemsep=0.15ex,topsep=0.2ex]
\item[5.2-1] 平坦陈-帕顿线丛的第一陈类为零，这是它能同时充当 $I$ 型 $B$-膜与 $J$、$K$ 型 $A$-膜的陈-帕顿丛的前提：$B$-膜只要求陈-帕顿丛是全纯丛，而 $A$-膜还要求它（在 K-理论意义下）满足适当的量子化条件. 对复环面 $\FF$，这样的平坦线丛由沿 $\FF$ 的一维闭链的和乐唯一指定（§5.3），故总是存在；自旋结构的微妙之处在于把膜提升为 $A$-膜时需选定 $\FF$ 上的一个 $\mathrm{Spin}^c$（或自旋）结构，而 $\FF$ 是复环面故可平行化，这种选择总存在，正文所考虑的 $\bf F$ 型膜因而不受阻碍. [来源：原文 §5.2 式 (5.1) 前后及关于 K-理论偏移与自旋结构的脚注；nLab “Kapustin-Witten TQFT”，\url{https://ncatlab.org/nlab/show/Kapustin-Witten+TQFT}（Last revised 2026-09-28；2026-11-19 访问）；2026-11-19 核对.]
\item[5.2-2] 原文此处两次写 $\L\to \L\otimes\O(\qq')$，按上下文第二处应为 $\L\to \L\otimes\O(\qq'')$：式 (5.1) 中 $\Phi'$、$\Phi''$ 分别对应 §4.5.2 的两个 $\varphi$-不变修正 $\L\to\L(\qq')$ 与 $\L\to\L(\qq'')$（式 (4.45)），而 §5.2 下文明说作用为 $\L\to\L\otimes\O(\qq')$ 或 $\L\to\L\otimes\O(\qq'')$. 这一处原文在 p.57 亦如此印刷，属原文笔误，译文照录. [来源：原文 §5.2 式 (5.1) 上下文及式 (4.45)；原文 §5.2 关于 $\Pic(D)$ 的段落（原文 p.57）；2026-11-19 核对.]
\end{description}
\medskip\noindent\textbf{第 5.3 节\quad 平移本征丛}\par\nopagebreak
\begin{description}[leftmargin=5.8em,labelwidth=5em,labelsep=0.6em,itemsep=0.15ex,topsep=0.2ex]
\item[5.3-1] “平移本征丛”的几何含义是：$\cF$ 是复环面，平移群是另一个复环面 $\t\cF$（选取基点后二者同构）；线丛 $\L\to \cF$ 沿平移的拉回 $f^*\L$ 与 $\L$ 同构，但同构不唯一（差一个 $\C^*$ 因子），所以“本征值”是一维空间 $\C$ 的非零元素，而不是一个数. 第一陈类为零的全纯线丛总容许与全纯结构相容的平坦联络，因此 $\cF$ 上的平坦线丛全部是平移本征丛；反之 $c_1(\L)\neq0$ 时不存在这样的结构. 高秩情形由式 (5.4)–(5.7) 化归为“固定向量空间 $\otimes$ 一个线丛”，故不会产生本质新的对象. [来源：原文 §5.3 式 (5.4)–(5.7) 及上下文；nLab “Hitchin fibration”，\url{https://ncatlab.org/nlab/show/Hitchin+fibration}（Last revised 2023-11-01；2026-11-19 访问）；2026-11-19 核对.]
\end{description}
\medskip\noindent\textbf{第 5.4 节\quad 作为磁本征膜的 $\bf F$ 型膜}\par\nopagebreak
\begin{description}[leftmargin=5.8em,labelwidth=5em,labelsep=0.6em,itemsep=0.15ex,topsep=0.2ex]
\item[5.4-1] 这里的“本征值”是 $\oplus_i{\cal N}_i$ 中各一维空间；${\cal N}_i$ 依赖于权 $m_i$ 的排序（即不同 $\Phi_i$），故不同排序给出的本征值不同，这与 §4.5.5 中“$\qq_i$ 的不同排序给出不同赫克修正”一致. 第 7 节将计算这些一维空间，其值与对偶的电本征膜一侧的预言相符，正是磁本征膜成为几何朗兰兹对应物理实现的关键. 更高秩的膜不会给出新的磁本征膜，因为 $\FF$ 上的任何平移本征丛按式 (5.7) 都是秩 1 本征丛与一个固定向量空间的张量积. [来源：原文 §5.4 式 (5.8) 及上下文；原文 §7；nLab “geometric Langlands correspondence”，\url{https://ncatlab.org/nlab/show/geometric+Langlands+correspondence}（2026-11-19 访问）；2026-11-19 核对.]
\end{description}
\medskip\noindent\textbf{第 6 节\quad 行列式线丛}\par\nopagebreak
\begin{description}[leftmargin=5.8em,labelwidth=5em,labelsep=0.6em,itemsep=0.15ex,topsep=0.2ex]
\item[6-1] 二维手征费米子在背景场中的路径积分在固定背景下是一个数，但它随背景（联络或复结构）变化时并不给出函数，而给出行列式线丛的一个截面；截面上的自然度量与自然联络就是 Quillen 度量与 Quillen 联络，其曲率即费米子行列式的反常。亏格 $g$ 的黎曼曲面上，狄拉克算符族的行列式（相差一个同构）正是 theta 函数，其零除子给出 theta 除子；时空维数为 $8k+2$ 时行列式线丛还有典范平方根，即 Pfaffian 线丛（黎曼曲面正属 $k=0$ 情形）。[来源：Quillen 1985（Funct. Anal. Appl. 19，31–34，本文文献 [41]）；原文 §6；nLab "determinant line bundle"，\url{https://ncatlab.org/nlab/show/determinant+line+bundle}（Last revised 2026-01-30；2026-11-19 访问）之 "Quillen's determinant line bundle"、"Pfaffian line bundle"、"Relation to theta function" 三节。]
\end{description}
\medskip\noindent\textbf{第 6.1 节\quad 线丛的配对}\par\nopagebreak
\begin{description}[leftmargin=5.8em,labelwidth=5em,labelsep=0.6em,itemsep=0.15ex,topsep=0.2ex]
\item[6.1-1] 注意式（\ref{bilpa}）右端是一维复向量空间（复线），不是复数：取行列式线即把上同调群取最高外幂，分式记号 $V/W$ 表示 $V\otimes W^{-1}$。因此“双线性”“对称”都只能理解为典范同构意义下的，这也解释了下面为何要强调 $\cong$ 可在族中同时给出。这一配对即 Deligne 的行列式—上同调配对（determinant of cohomology 配对）。[来源：Deligne 1987（Contemp. Math. 67，93–177，本文文献 [40]）；原文 §6.1；nLab "determinant line bundle"（2026-11-19 访问）。]
\end{description}
\medskip\noindent\textbf{第 6.2 节\quad 解释}\par\nopagebreak
\begin{description}[leftmargin=5.8em,labelwidth=5em,labelsep=0.6em,itemsep=0.15ex,topsep=0.2ex]
\item[6.2-1] 局部地看，若把万有线丛换为 $W\otimes R$（$R$ 为 $\Pic_c(D)$ 上的线丛），式（\ref{woffo}）表明 $\N$ 的转移函数只在一个 $d$ 次上闭链（2-cocycle）意义下良定义；这正是“被 gerbe 的 $d$ 次幂扭曲的扭曲线丛”。$d$ 是 $\M$ 的次数，即总电荷，所以 §7 计算的赫克“本征值”只能是这种扭曲意义下的值；要得到真正的数，须取 $d=0$ 或先固定一个平凡化。[来源：原文 §6.2–§6.3；[KW]（arXiv:hep-th/0604151）§7；nLab "gerbe"，\url{https://ncatlab.org/nlab/show/gerbe}（2026-11-19 访问）。]
\end{description}
\medskip\noindent\textbf{第 6.3 节\quad Gerbe}\par\nopagebreak
\begin{description}[leftmargin=5.8em,labelwidth=5em,labelsep=0.6em,itemsep=0.15ex,topsep=0.2ex]
\item[6.3-1] $C$ 上的自旋结构可视为典范丛 $K$ 的平方根（theta characteristic）；任意两个相差一个二阶线丛，这正对应 $\Z_2$ gerbe 的“带”。对一族狄拉克算符（时空维数 $8k+2$ 时），行列式线丛有典范平方根（Pfaffian 线丛），但把它识别为某个线丛需要先选自旋结构，故平凡化非典范；这与脚注 \ref{brief} 中“希钦纤维的等同依赖 $K^{1/2}$ 的选择”是同一现象。[来源：原文 §6.3 与 §6.5 脚注；nLab "determinant line bundle" 之 "Pfaffian line bundle" 与 "square roots" 表（canonical bundle 的平方根 = theta characteristic = spin structure）（2026-11-19 访问）。]
\item[6.3-2] 此 gerbe 的“带”（band）是中心 $\Z_n$：$^L\neg G$ 的万有丛在整体上不存在，因为中心在表示纤维上非平凡地作用；任意两个局部万有丛相差一个 $n$ 阶线丛的扭曲。故其障碍类落在 $H^2(\M(^L\neg G,C),\Z_n)$，非平凡性即“离散扭转”（discrete torsion）；[KW] 第 7 节正是用这一点研究 $^L\neg G$-丛的扭曲向量丛。[来源：[KW]（arXiv:hep-th/0604151）§7；原文 §6.3；nLab "gerbe" 之 "Banded gerbes" 节（2026-11-19 访问）。]
\item[6.3-3] 这里的 $B$ 即物理中的 $B$ 场（2-形式势）：改变平凡化时 $B\to B+F$ 是“大规范变换”，$H=dB$ 是规范不变的三形式场强；$B/2\pi$ 的周期取整为通量量子化条件。平坦 gerbe（$H=0$）上的平坦扭曲线丛，就是低能有效理论中“平坦 $B$ 场”的几何对应物。[来源：原文 §6.3；[KW]（arXiv:hep-th/0604151）§7；nLab "gerbe"（2026-11-19 访问）及其子条目 "gerbe (in differential geometry)"。]
\end{description}
\medskip\noindent\textbf{第 6.4 节\quad 对偶性}\par\nopagebreak
\begin{description}[leftmargin=5.8em,labelwidth=5em,labelsep=0.6em,itemsep=0.15ex,topsep=0.2ex]
\item[6.4-1] 式（\ref{pliko}）是 Quillen 度量的曲率公式：对固定的光滑复线丛取一族酉联络（即一族 $\bar\partial$-算符），其行列式线丛带自然厄米度量与自然酉联络，曲率即式（\ref{pliko}）。$\langle{\cal L},{\cal M}\rangle$ 的曲率由式（\ref{zunko}）中四个行列式线丛的曲率按分式的 $\pm$ 取线性组合得到，即式（\ref{likko}）。当 ${\cal L},{\cal M}$ 限制为平坦联络、即 $\Jac(D)$ 上的点时，这个 2-形式是 $\Jac(D)\times\Jac(D)$ 的辛（Kähler）形式（相差归一化），故具此曲率的酉线丛必为庞加莱线丛。[来源：原文 §6.4；Quillen 1985（Funct. Anal. Appl. 19，31–34，本文文献 [41]）；2026-11-19 核对。]
\end{description}
\medskip\noindent\textbf{第 6.5 节\quad 再论希钦纤维的对偶性}\par\nopagebreak
\begin{description}[leftmargin=5.8em,labelwidth=5em,labelsep=0.6em,itemsep=0.15ex,topsep=0.2ex]
\item[6.5-1] 在谱覆盖 $\psi:D\to C$ 下，$U(N)$ 的希钦纤维（选定 $K^{1/2}$ 后）可等同于 $D$ 的雅可比簇（同构依赖 $K^{1/2}$ 的选择）；§6.4 的配对说明 $\Jac(D)$ 与其对偶阿贝尔簇可由庞加莱线丛等同，谱覆盖再把这一等同翻译成希钦纤维的对偶性。注意这里的纤维是阿贝尔簇（可带扭化），并非字面意义上的环面。[来源：原文 §6.5；Hitchin 1987（Proc. London Math. Soc. 55，59–126，本文文献 [13]）；2026-11-19 核对。]
\item[6.5-2] 式（\ref{unofto}）说 $\det H^*(C,E)$ 作为 $E$ 的函数只通过 $\det E$ 分解：右端秩为 1，只含 $\det E$ 的信息，而左端本可依赖 $E$ 的高阶陈类，故这一恒等式并不显然。证明用线子丛 ${\cal M}\subset E$、正合列（\ref{gonno}）与相应的长正合列，把秩 $r$ 的情形化到秩 $r-1$；§6.5 正是用它证明：固定 ${\rm Nm}({\cal L})$ 时 $\langle{\cal L},\psi^*{\cal M}_0\rangle$ 不随 ${\cal L}$ 变化。[来源：原文 §6.5；2026-11-19 核对。]
\end{description}
\medskip\noindent\textbf{第 7 节\quad 计算赫克本征值}\par\nopagebreak
\begin{description}[leftmargin=5.8em,labelwidth=5em,labelsep=0.6em,itemsep=0.15ex,topsep=0.2ex]
\item[7-1] 对 $U(1)$ 来说中心与基本群都是 $\Z$：电荷 $d\in\Z$ 无界，磁通量（即 $c_1$）也可取任意整数，故 't Hooft 算符能把 $c_1({\cal L})$ 平移任意整数，把膜移出 $\MH$。对单非阿贝尔群，中心有限（如 $SU(N)$ 的中心为 $\Z_N$），可平移的 $c_1$ 只能是 $N$ 的倍数，于是磁通量成为离散不变量 $\xi$（§2.6.6）；这正是阿贝尔情形必须额外保留 $U(1)$ 规范多重态的根本原因。[来源：原文 §7、§7.1.2、§2.6.6；2026-11-19 核对。]
\end{description}
\medskip\noindent\textbf{第 7.1.2 节\quad 阿贝尔情形}\par\nopagebreak
\begin{description}[leftmargin=5.8em,labelwidth=5em,labelsep=0.6em,itemsep=0.15ex,topsep=0.2ex]
\item[7.1.2-1] 可验算其归一化：由 $\int_C\omega=1$ 与 $F=2\pi d\,\omega$ 得 $\int_C F=2\pi d$，与 $c_1({\cal L})=d$ 的 Chern–Weil 表示一致；又 $\phi_0=2\pi y\,d$ 使 $\partial_y\phi_0=2\pi d$，正是后文“$\phi_0$ 的法向导数等于 $2\pi d$”的边界条件。故 $d\not=0$ 时此解给出带磁荷 $d$ 的膜，而非 $\MH$ 上的零维膜——低能描述必须含 $U(1)$ 向量多重态。[来源：原文 §7.1.2；[KW]（arXiv:hep-th/0604151）式（3.29）取 $t=1$；2026-11-19 核对。]
\end{description}
\medskip\noindent\textbf{第 7.2.2 节\quad 非阿贝尔推广}\par\nopagebreak
\begin{description}[leftmargin=5.8em,labelwidth=5em,labelsep=0.6em,itemsep=0.15ex,topsep=0.2ex]
\item[7.2.2-1] $k+rN$ 是权 $^Lw(k;r)=(r+1,\dots,r+1,r,\dots,r)$ 之和，即表示 $\wedge^k{\cal V}\otimes(\det{\cal V})^r$ 的权和；'t Hooft 算符按此权平移 $c_1(E)$，故 $\Delta$ 就是移位量。取 $^L\neg G=SU(N)$（对应 $G=SU(N)/\Z_N$）时权在 $\Z_N$ 上的和为零，$\Delta$ 只能是 $N$ 的倍数，而 $PSU(N)$ 描述中的 $c_1$ 本就只定到模 $N$，于是不必额外保留 $U(1)$ 多重态——这正是本节末改用 $SU(N)$ 的理由。[来源：原文 §7.2.2；2026-11-19 核对。]
\end{description}
\endgroup
